\documentclass[12pt, letterpaper]{article}

\usepackage[T1]{fontenc}
\usepackage[latin1]{inputenc}
\usepackage[spanish]{babel}
\usepackage{geometry}
\geometry{margin=2.5cm}
\usepackage{setspace}
\onehalfspacing
\usepackage{parskip}
\setlength{\parskip}{8pt}
\usepackage{hyperref}
\hypersetup{colorlinks=true, linkcolor=blue, citecolor=blue, urlcolor=blue}

\title{Modelaci\'on de la Capa Mixta Oce\'anica en el Pac\'ifico Colombiano:
	Respuesta T\'ermica y Precipitaci\'on bajo Escenarios de Calentamiento Global}
\author{Jhon Garc\'ia Barrera \\
	F\'isica del Clima y el Cambio Clim\'atico \\
	Universidad del Quind\'io}
\date{\today}

\begin{document}
	
	\maketitle
	
	Colombia es uno de los pa\'ises con mayor disponibilidad h\'idrica del planeta, y esa riqueza
	depende de manera cr\'itica del comportamiento de la precipitaci\'on sobre su vertiente
	Pac\'ifica y sus Andes occidentales. Esa precipitaci\'on est\'a controlada en gran medida
	por la temperatura superficial del mar (SST) en el Pac\'ifico colombiano (aproximadamente
	0--8$^\circ$N, 75--82$^\circ$W), una regi\'on donde la Zona de Convergencia Intertropical
	(ITCZ) y el Chorro del Choc\'o (la corriente de bajo nivel que transporta la humedad desde
	el oc\'eano hacia el interior del pa\'is) interact\'uan estrechamente con las condiciones
	oce\'anicas de superficie [1]. Esta conexi\'on no es nueva: Poveda et al. [2] demostraron
	que las anomal\'ias de SST asociadas al fen\'omeno El Ni\~no-Oscilaci\'on del Sur (ENSO)
	explican hasta el 50\,\% de la variabilidad hidrol\'ogica colombiana, y Poveda et al. [3]
	extendieron ese resultado mostrando influencia adicional de la Oscilaci\'on Decadal del
	Pac\'ifico y del Atl\'antico tropical. M\'as recientemente, C\'ordoba-Machado et al. [4]
	discriminaron que la precipitaci\'on en la Regi\'on Andina responde principalmente al
	Pac\'ifico central (El Ni\~no Modoki), mientras que la Regi\'on Pac\'ifica responde al
	Pac\'ifico oriental, donde la SST local act\'ua como modulador directo de la inestabilidad
	convectiva.
	
	El problema que este proyecto aborda es el siguiente: si bien la conexi\'on entre SST y
	precipitaci\'on en Colombia est\'a documentada en el contexto de la variabilidad interanual,
	no existe una cuantificaci\'on regional de c\'omo responder\'ia la SST del Pac\'ifico
	colombiano a un forzamiento radiativo sostenido (como el que impone el aumento de gases
	de efecto invernadero) y cu\'ales ser\'ian las consecuencias hidrol\'ogicas de esa
	respuesta. Las proyecciones globales de CMIP6 [5] ofrecen informaci\'on sobre la SST futura
	a escala de cuenca, pero su resoluci\'on espacial es insuficiente para capturar la din\'amica
	local de una regi\'on tan compleja como el litoral Pac\'ifico colombiano, donde la
	orograf\'ia andina, el Chorro del Choc\'o y la posici\'on de la ITCZ interact\'uan a escalas
	de pocos grados de latitud. Adicionalmente, los modelos globales muestran una sensibilidad
	clim\'atica m\'as alta en CMIP6 que en generaciones anteriores precisamente por cambios en
	las retroalimentaciones de nubes sobre el oc\'eano tropical [6], lo que hace m\'as urgente
	disponer de estimaciones regionales robustas.
	
	La capa mixta oce\'anica (CML) es la capa superficial del oc\'eano que interact\'ua
	directamente con la atm\'osfera. Su temperatura evoluciona en respuesta al balance entre el
	calentamiento solar, el enfriamiento por radiaci\'on infrarroja, los flujos turbulentos de
	calor latente y sensible, y el enfriamiento por entrainment de aguas m\'as fr\'ias desde
	abajo [7]. En el Pac\'ifico tropical oriental, donde la CML es delgada (entre 15 y 50 m
	seg\'un la climatolog\'ia global de de Boyer Mont\'egut et al. [8], actualizada con datos
	Argo recientemente [9]) perturbaciones peque\~nas en el forzamiento superficial producen
	cambios de SST amplificados, lo que hace a esta regi\'on especialmente sensible al
	calentamiento global. El modelo unidimensional de Price, Weller y Pinkel (PWP) [10] es el
	est\'andar de referencia para simular esta din\'amica: incorpora ajuste convectivo ante
	inestabilidad est\'atica y mezcla por inestabilidad de cizalladura mediante un n\'umero de
	Richardson de bulto. Los flujos turbulentos en la interfaz aire-mar se calculan mediante el
	algoritmo COARE, desarrollado por Fairall et al. [11] y extendido en su versi\'on 3.0 [12],
	que fue dise\~nado espec\'ificamente para condiciones de viento d\'ebil y convecci\'on
	intensa, exactamente las condiciones prevalecientes en el Pac\'ifico colombiano.
	
	El marco te\'orico para conectar los cambios de SST con cambios de precipitaci\'on est\'a
	bien establecido. Xie et al. [13] demostraron que bajo calentamiento global las lluvias
	aumentan donde el calentamiento local de la SST supera la media tropical, y disminuyen donde
	queda por debajo. Xie [14] extendi\'o ese marco mostrando que el patr\'on espacial de
	calentamiento de SST es el principal determinante de los cambios regionales de precipitaci\'on
	en el tr\'opico, incluso sobre el efecto termodin\'amico del aumento uniforme de humedad.
	Johnson y Xie [15] precisaron que el umbral de SST para la convecci\'on profunda se desplaza
	con la media tropical, de modo que la respuesta de precipitaci\'on a una perturbaci\'on de
	SST es predecible a primer orden. Sun et. al. [16] a\~nadieron que esta sensibilidad tiene
	una marcada dependencia estacional en el Pac\'ifico ecuatorial oriental, lo que subraya la
	necesidad de resolver el ciclo anual correctamente. Por su parte, Vecchi y Soden [17]
	mostraron que el debilitamiento de la circulaci\'on tropical bajo calentamiento global altera
	la posici\'on de la ITCZ y los patrones de convergencia de humedad, alterando la posici\'on media de la ITCZ en funci\'on del balance energ\'etico interhemisf\'erico [18], lo que tiene implicaciones directas para la precipitaci\'on colombiana [1, 3]
	
El camino que se propone seguir tiene tres etapas. La primera consiste en implementar en
Python el modelo PWP forzado con el rean\'alisis ERA5 [19], que provee variables atmosf\'ericas
horarias a 0.25$^\circ$ de resoluci\'on para el per\'iodo 1980--2023. La segunda etapa es
la calibraci\'on y validaci\'on del modelo: la SST simulada se comparar\'a contra el producto
OISST v2.1 de NOAA [20] (SST diaria observada basada en sat\'elites AVHRR y boyas in
situ) mediante RMSE, sesgo estacional y correlaci\'on de Pearson, verificando adem\'as que
el modelo reproduzca los eventos ENSO hist\'oricos documentados para el Pac\'ifico colombiano
por Cer\'on et al. [21] y Morales-Mar\'in et al. [22]. La tercera etapa es experimental: una
vez validado, el modelo se forzar\'a con perturbaciones de temperatura del aire de
$+1.5^\circ$C, $+2^\circ$C y $+4^\circ$C, coherentes con los escenarios SSP1-2.6, SSP2-4.5
	y SSP5-8.5 de CMIP6 [5]. La diferencia de SST entre cada escenario y el control se
	traducir\'a en cambios de precipitaci\'on usando una relaci\'on emp\'irica SST--precipitaci\'on
	calibrada con datos CHIRPS v2 [23] para la regi\'on costera y andina adyacente, siguiendo
	el marco de Xie et al. [13] y el balance de calor de la capa mixta descrito por Swenson y
	Hansen [7].
	
	El aporte central del proyecto es doble. Por un lado, implementa f\'isica local expl\'icita
	de la interfaz aire-mar con resoluci\'on espacial superior a la de los modelos globales,
	lo que permite capturar la respuesta de la SST en una regi\'on con din\'amica propia. Por
	otro lado, traduce esa respuesta en estimaciones de impacto hidrol\'ogico regional directamente
	relevantes para Colombia. No existe precedente publicado de este enfoque para el Pac\'ifico
	colombiano, lo que constituye la originalidad cient\'ifica del trabajo.
	
	\bigskip
	\noindent\textbf{Referencias}
	
	\smallskip
	\noindent [1] Poveda, G. (2004). La hidroclimatolog\'ia de Colombia: una s\'intesis desde la
	escala interdecadal hasta la escala diurna. \textit{Revista de la Academia Colombiana de
		Ciencias Exactas, F\'isicas y Naturales}, 28(107), 201--222.
	
	\noindent [2] Poveda, G., Jaramillo, A., Gil, M.\,M., Quiceno, N., y Mantilla, R.\,I. (2001).
	Seasonality in ENSO-related precipitation, river discharges, soil moisture, and vegetation
	index in Colombia. \textit{Water Resources Research}, 37(8), 2169--2178.
	\url{https://doi.org/10.1029/2000WR900395}
	
	\noindent [3] Poveda, G., \'Alvarez, D.\,M., y Rueda, \'O.\,A. (2011). Hydro-climatic
	variability over the Andes of Colombia associated with ENSO: a review of climatic processes
	and their impact on one of the Earth's most important biodiversity hotspots.
	\textit{Climate Dynamics}, 36(11), 2233--2249.
	\url{https://doi.org/10.1007/s00382-010-0931-y}
	
	\noindent [4] C\'ordoba-Machado, S., Palomino-Lemus, R., G\'amiz-Fortis, S.\,R.,
	Castro-D\'iez, Y., y Esteban-Parra, M.\,J. (2015). Assessing the impact of El Ni\~no Modoki
	on seasonal precipitation in Colombia. \textit{Global and Planetary Change}, 124, 41--61.
	\url{https://doi.org/10.1016/j.gloplacha.2014.11.003}
	
	\noindent [5] Eyring, V., Bony, S., Meehl, G.\,A., Senior, C.\,A., Stevens, B., Stouffer,
	R.\,J., y Taylor, K.\,E. (2016). Overview of the Coupled Model Intercomparison Project
	Phase 6 (CMIP6) experimental design and organization. \textit{Geoscientific Model
		Development}, 9(5), 1937--1958.
	\url{https://doi.org/10.5194/gmd-9-1937-2016}
	
	\noindent [6] Zelinka, M.\,D., Myers, T.\,A., McCoy, D.\,T., Po-Chedley, S., Caldwell,
	P.\,M., Ceppi, P., Klein, S.\,A., y Taylor, K.\,E. (2020). Causes of higher climate
	sensitivity in CMIP6 models. \textit{Geophysical Research Letters}, 47(1), e2019GL085782.
	\url{https://doi.org/10.1029/2019GL085782}
	
	\noindent [7] Swenson, M.\,S., y Hansen, D.\,V. (1999). Tropical Pacific Ocean mixed layer
	heat budget: The Pacific cold tongue. \textit{Journal of Physical Oceanography}, 29(1),
	69--81.
	\url{https://doi.org/10.1175/1520-0485(1999)029<0069:TPOMLH>2.0.CO;2}
	
	\noindent [8] de Boyer Mont\'egut, C., Madec, G., Fischer, A.\,S., Lazar, A., e Iudicone,
	D. (2004). Mixed layer depth over the global ocean: An examination of profile data and a
	profile-based climatology. \textit{Journal of Geophysical Research: Oceans}, 109(C12),
	C12003.
	\url{https://doi.org/10.1029/2004JC002378}
	
	\noindent [9] de Boyer Mont\'egut, C. (2024). Mixed layer depth over the global ocean: a
	climatology computed with a density threshold criterion of 0.03 kg/m$^3$ from 5 m depth
	value. SEANOE.
	\url{https://doi.org/10.17882/98226}
	
	\noindent [10] Price, J.\,F., Weller, R.\,A., y Pinkel, R. (1986). Diurnal cycling:
	Observations and models of the upper ocean response to diurnal heating, cooling, and wind
	mixing. \textit{Journal of Geophysical Research: Oceans}, 91(C7), 8411--8427.
	\url{https://doi.org/10.1029/JC091iC07p08411}
	
	\noindent [11] Fairall, C.\,W., Bradley, E.\,F., Rogers, D.\,P., Edson, J.\,B., y Young,
	G.\,S. (1996). Bulk parameterization of air-sea fluxes for Tropical Ocean Global Atmosphere
	Coupled Ocean-Atmosphere Response Experiment. \textit{Journal of Geophysical Research:
		Oceans}, 101(C2), 3747--3764.
	\url{https://doi.org/10.1029/95JC03205}
	
	\noindent [12] Fairall, C.\,W., Bradley, E.\,F., Hare, J.\,E., Grachev, A.\,A., y Edson,
	J.\,B. (2003). Bulk parameterization of air-sea fluxes: Updates and verification for the
	COARE algorithm. \textit{Journal of Climate}, 16(4), 571--591.
	\url{https://doi.org/10.1175/1520-0442(2003)016<0571:BPOASF>2.0.CO;2}
	
	\noindent [13] Xie, S.-P., Deser, C., Vecchi, G.\,A., Ma, J., Teng, H., y Wittenberg,
	A.\,T. (2010). Global warming pattern formation: Sea surface temperature and rainfall.
	\textit{Journal of Climate}, 23(4), 966--986.
	\url{https://doi.org/10.1175/2009JCLI3329.1}
	
	\noindent [14] Xie, S.-P. (2020). Ocean warming pattern effect on global and regional
	climate change. \textit{AGU Advances}, 1(1), e2019AV000130.
	\url{https://doi.org/10.1029/2019AV000130}
	
	\noindent [15] Johnson, N.\,C., y Xie, S.-P. (2010). Changes in the sea surface temperature
	threshold for tropical convection. \textit{Nature Geoscience}, 3, 842--845.
	\url{https://doi.org/10.1038/ngeo1008}
	
	\noindent [16] Sun, N., Zhou, T., Chen, X., Endo, H., Kitoh, A., y Wu, B. (2020).
	Amplified tropical Pacific rainfall variability related to background SST warming.
	\textit{Climate Dynamics}, 54, 2387--2402.
	\url{https://doi.org/10.1007/s00382-020-05119-3}
	
	\noindent [17] Vecchi, G.\,A., y Soden, B.\,J. (2007). Global warming and the weakening of
	the tropical circulation. \textit{Journal of Climate}, 20(17), 4316--4340.
	\url{https://doi.org/10.1175/JCLI4258.1}
	
	\noindent [18] Schneider, T., Bischoff, T., y Haug, G.\,H. (2014). Migrations and dynamics
	of the intertropical convergence zone. \textit{Nature}, 513, 45--53.
	\url{https://doi.org/10.1038/nature13636}
	
	\noindent [19] Hersbach, H., Bell, B., Berrisford, P., y colaboradores (2020). The ERA5
	global reanalysis. \textit{Quarterly Journal of the Royal Meteorological Society}, 146(730),
	1999--2049.
	\url{https://doi.org/10.1002/qj.3803}
	
	\noindent [20] Huang, B., Liu, C., Banzon, V., Freeman, E., Graham, G., Hankins, B., Smith,
	T., y Zhang, H.-M. (2021). Improvements of the Daily Optimum Interpolation Sea Surface
	Temperature (DOISST) version 2.1. \textit{Journal of Climate}, 34(8), 2923--2939.
	\url{https://doi.org/10.1175/JCLI-D-20-0166.1}
	
	\noindent [21] Cer\'on, W.\,L., Kayano, M.\,T., Andreoli, R.\,V., Canchala, T.,
	Carvajal-Escobar, Y., y Alfonso-Morales, W. (2021). Rainfall variability in
	southwestern Colombia: changes in ENSO-related features.
	\textit{Pure and Applied Geophysics}, 178(3), 1087--1103.
	\url{https://doi.org/10.1007/s00024-021-02673-7}
	
	\noindent [22] Morales-Mar\'in, L.\,A., Rodr\'iguez, E.\,A., y Jaramillo, F. (2026).
	Water resources challenges in Colombia: hydrological science solutions for a
	sustainable future. \textit{Hydrological Sciences Journal}, 1--20.
	\url{https://doi.org/10.1080/02626667.2026.2631143}
	
	\noindent [23] Funk, C., Peterson, P., Landsfeld, M., Pedreros, D., Verdin, J., Shukla, S.,
	Husak, G., Rowland, J., Harrison, L., Hoell, A., y Michaelsen, J. (2015). The climate
	hazards infrared precipitation with stations---a new environmental record for monitoring
	extremes. \textit{Scientific Data}, 2, 150066.
	\url{https://doi.org/10.1038/sdata.2015.66}
	
\end{document}